\chapter{Conclusiones y Trabajo Futuro}
\label{cap:conclusiones}

\chapterquote{Sólo los optimistas se suicidan, los optimistas que ya no pueden serlo. Los demás, no teniendo ninguna razón para vivir, tampoco la tienen para morir.}{El inconveniente de haber nacido - E.Cioran}

\section{Conclusiones}

En este trabajo se plantearon dos objetivos principales:

\begin{itemize}
    \item \textbf{Incorporar funcionalidades de realidad aumentada (RA)} dentro de la plataforma \textit{OBS Studio}.
    \item \textbf{Aplicar dicha RA en concursos de programación en directo} mediante el desarrollo de un plugin nativo.
\end{itemize}

Ambos objetivos se han cumplido satisfactoriamente.

El resultado es un plugin funcional que extiende \textit{OBS Studio} con cuatro modos de operación diferenciados. En el modo de modelo 3D, el sistema carga geometría en múltiples formatos mediante \textbf{ASSIMP}, estima la pose del marcador \textbf{ArUco} en tiempo real usando el algoritmo PnP de \textbf{OpenCV} y renderiza el modelo anclado al marcador con la orientación y posición correctas. En el modo cronómetro, ese mismo modelo actúa como reloj analógico de cuenta atrás cuyas manecillas se actualizan y se sincronizan periódicamente con el tiempo oficial conectado a la API de \textbf{DOMjudge} \citep{DOMjudge}. En el modo \textit{Scoreboard}, el sistema proyecta la clasificación general del concurso sobre el vídeo en tiempo real, con los datos de todos los equipos organizados. En el modo \textit{Team Info}, la cámara detecta automáticamente el marcador de un equipo específico y muestra únicamente sus datos, combinando la detección de marcadores con un fichero JSON de mapeo entre equipos que el usuario prepara antes de la retransmisión.

Más allá de los modos de visualización, el trabajo ha resuelto varios problemas técnicos. El sistema de coordenadas que separa OpenCV de OBS requirió una transformación explícita de convenios mediante rotaciones. El suavizado del overlay mediante filtro, elimina gran parte del temblor perceptible que producen el ruido de cámara y la compresión de vídeo. La arquitectura de dos hilos, con un hilo secundario dedicado exclusivamente a las peticiones HTTP, garantiza que la red nunca bloquee el renderizado. El proceso de calibración de cámara, implementado como herramienta independiente, genera el fichero de parámetros que el plugin carga al arrancar y que permite corregir la distorsión óptica de cada objetivo.

El sistema de conversión de formatos de vídeo cubre los siete formatos que OBS puede entregar (I420, NV12, YUY2, UYVY, BGRA, RGBA, Y800), haciendo que la detección funcione independientemente de la fuente de cámara conectada, mientras cumpla con esos requisitos. La interfaz de usuario integrada en OBS permite al usuario ajustar todos los parámetros en directo, en cualquier momento, sin reiniciar la escena. Esto hace el sistema manejable en un entorno de producción real.

En conjunto, el plugin demuestra que es posible añadir un sistema completo de realidad aumentada orientada a datos en tiempo real a \textit{OBS Studio} manteniendo la arquitectura de plugin nativo, sin latencia adicional por procesos externos y sin modificar el núcleo del software.

\section{Limitaciones}

El plugin ha sido desarrollado y probado en un entorno Windows con la cámara integrada del ASUS TUF Gaming A16. Esta circunstancia implica varias limitaciones que conviene señalar explícitamente. El sistema no ha sido validado en otros entornos como Linux o macOS, plataformas en las que OBS Studio también funciona, pero cuya API y gestión de recursos difieren. Además la calidad de la detección depende directamente de las condiciones de iluminación y de la resolución de impresión de los marcadores. La calibración de la cámara debe repetirse para cada objetivo diferente, lo que añade un paso de configuración que puede resultar poco práctico en entornos donde se cambia de cámara con frecuencia. Por último, el sistema de conectividad está diseñado exclusivamente para la API REST de \textbf{DOMjudge} \citep{DOMjudge}, por lo que no puede conectarse a otros sistemas de gestión de concursos de programación sin modificar el código realizado en este plugin.

\section{Trabajo Futuro}

Los resultados obtenidos abren varias líneas de continuación natural.

La más relevante desde el punto de vista técnico es la incorporación de \textbf{seguimiento sin marcadores} (\textit{markerless tracking}). Las técnicas actuales de SLAM visual \citep{MurArtal2015} permitirían anclar los elementos de RA a superficies del entorno real. Esto no tendría la necesidad de imprimir y colocar códigos ArUco en diferentes elementos. Por ello haría al sistema más flexible y aplicable a escenarios donde no se puede controlar el espacio físico.

Una segunda línea es la \textbf{compatibilidad con otros sistemas de concursos en directo}. Actualmente el proyecto está adaptado a la API de DOMjudge, por lo que añadir soporte para otras plataformas ampliaría el alcance del plugin a una sección mucho mayor de la comunidad de concursos de programación.

En el plano del rendimiento, sería valioso realizar una \textbf{evaluación sistemática del impacto} del plugin sobre la tasa de fotogramas y el uso de CPU y GPU en diferentes configuraciones. Esto establecería umbrales de uso recomendados y posibles optimizaciones para diferentes equipos.

Finalmente, añadir un \textbf{panel de control remoto} accesible desde una interfaz web o una aplicación móvil. Esto permitiría al director de realización gestionar los modos y la visibilidad de los elementos de RA desde un dispositivo secundario. Por ello no tendría la necesidad de tocar el equipo de retransmisión durante la emisión. Esto mejoraría significativamente el flujo de trabajo en producciones con más de una persona en el equipo técnico.
